\documentclass[11pt]{article}
\usepackage[margin=1in]{geometry}

\title{A Minimal Manuscript for Exercising the Paper Review Contract}
\author{Paper Review Bench Maintainers}
\date{}

\begin{document}
\maketitle

\begin{abstract}
This short synthetic manuscript is a complete, self-contained input for an
end-to-end paper-review smoke test. It states one elementary claim, gives its
derivation, and reports a small consistency check. The scientific content is
deliberately modest: the task is to review the manuscript as submitted, not to
infer an unstated result or search for outside literature.
\end{abstract}

\section{Claim}\label{sec:claim}
For every real number $x$, the quantity $(x-1)^2$ is non-negative. Equivalently,
\begin{equation}\label{eq:bound}
x^2 + 1 \geq 2x.
\end{equation}
We use this identity as a small example of a claim whose statement, argument,
and stated limitations can all be inspected directly in the supplied source.

\section{Derivation}\label{sec:derivation}
Expanding the square gives
\[
(x-1)^2 = x^2 - 2x + 1.
\]
A square of a real number is non-negative, so rearranging the displayed
identity proves Equation~\ref{eq:bound}. No additional assumptions are used.

\section{Consistency Check}\label{sec:check}
Table~\ref{tab:values} evaluates both sides for three representative inputs.
The values are consistent with the claimed inequality.

\begin{table}[h]
\centering
\begin{tabular}{c|c|c}
$x$ & $x^2 + 1$ & $2x$ \\
\hline
0 & 1 & 0 \\
1 & 2 & 2 \\
2 & 5 & 4 \\
\end{tabular}
\caption{A direct finite check of Equation~\ref{eq:bound}.}
\label{tab:values}
\end{table}

\section{Scope}\label{sec:scope}
This manuscript establishes only the elementary inequality in
Equation~\ref{eq:bound}. It makes no novelty claim, empirical generalization,
or claim about a production system. A useful review can therefore assess
whether the claim is stated precisely, the derivation is valid, the check is
appropriately interpreted, and the scope is communicated clearly.

\end{document}
